\section{TESTES DO NIST 800-22}
\subsection{Teste de Frequência}%PASSO 1 - Explicacao do resultado do (Teste de Frequência (monobit))
            O teste de frequência (monobit) é o teste mais básico da bateria do NIST.\@
            É o pilar para a execução dos demais testes. Se um gerador não for capaz de equilibrar a quantidade de bits '1' e '0' em uma sequência, ele falha nesse teste fundamental, indicando que a sequência não é aleatória o suficiente para aplicações criptográficas.
            Verificação do resultado do teste de frequência (monobit) para ambos os geradores, destacando a falha do java.util.Random e o sucesso do SecureRandom.

            \begin{figure}[htb]
                    \centering
                    \caption{Resultado do Teste de Frequência (Monobit) para SecureRandom}\label{fig:01_segure}
                    \includegraphics[width=0.8\textwidth]{NIST/01/01_segure.png}
                    \legend{Fonte: Elaborado pelo autor (2026)}
            \end{figure}

            Cada valor apresentado no resultado do teste de frequência (monobit) representa o $P\text{-value}$ calculado para a sequência de bits gerada. 
            Ele verifica se a quantidade de bits '1' e '0' está equilibrada, como seria esperado em uma sequência verdadeiramente aleatória.
            
            $P\text{-value}$ proximo de 1 indica que possui uma distribuição equilibrada, como é o caso do valor 0.985639.
            
            $P\text{-value}$ proximo de 0 indica que possui uma distribuição desequilibrada, como é o caso do valor 0.022021.
            
            $P\text{-value}$ menor que 0.01 é considerado uma falha no teste, indicando que a sequência não é aleatória o suficiente, como é o caso do valor 0.000128.

            \begin{figure}[htb]
                    \centering
                    \caption{Resultado final do Teste de Frequência (Monobit) para SecureRandom}\label{fig:01_final_segure}
                    \includegraphics[width=0.8\textwidth]{NIST/01/01_final_segure.png}
                    \legend{Fonte: Elaborado pelo autor (2026)}
            \end{figure}

            Os numeros apresentados no resultado final do teste de frequência (monobit). Da coluna C1 até a coluna C10 representam a distribuição das 100 amostras.
            Que foi dividida em 10 blocos de 10 amostras cada. O valor da coluna P representa a proporção de amostras que passaram no teste, ou seja, que tiveram um $P\text{-value}$ maior ou igual a 0.01. O valor da coluna P' representa o $P\text{-value}$ calculado para a proporção de amostras.
            Cada amostra tiveram um $P\text{-value}$ entre 0.0 e 1.0.
            Na coluna C1 nos é apresentado que obtivemos 10 amostras com $P\text{-value}$ maior ou igual a 0.01.
            Na coluna C10 nos é apresentado que obtivemos 12 amostras com $P\text{-value}$ entre 0.9 a 1.00.
            O $P\text{-value}$ 0.739918 apresentado é um teste Qui-Quadrado aplicado sobre os números das colunas C1 a C10. Fazendo a medição da uniformidade e a distribuição dos resultados.
            Se este valor for maior ou igual a 0.0001, o teste é considerado bem sucedido, indicando que os resultados estão distribuídos de maneira uniforme.
            Com o resultado de 0.739918, o SegureRandom passou no teste de frequência (monobit) com uma distribuição de $P\text{-value}$ que é considerada exelente. 
            Capaz de produzir sequências de bits que se aproximam muito da aleatoriedade ideal, o que é crucial para aplicações criptográficas seguras.
            Com uma proporsão de passagem de 99/100. Significa que 99 amostras obtiveram um $P\text{-value}$ maior ou igual a 0.01. 
            Para uma amostragem de 100, o minimo requisitado para passar no teste é de 96/100.

            \begin{figure}[htb]
                    \centering
                    \caption{Resultado do Teste de Frequência (Monobit) para Random}\label{fig:01_insegure}
                    \includegraphics[width=0.8\textwidth]{NIST/01/01_insegure.png}
                    \legend{Fonte: Elaborado pelo autor (2026)}
            \end{figure}
            \begin{figure}[htb]
                    \centering
                    \caption{Resultado final do Teste de Frequência (Monobit) para Random}\label{fig:01_final_insegure}
                    \includegraphics[width=0.8\textwidth]{NIST/01/01_final_insegure.png}
                    \legend{Fonte: Elaborado pelo autor (2026)}
            \end{figure}

            Por mais que olhando os valores apresentados no resultado do teste de frequência (monobit) para o Random,
            seja possível observar que a maioria dos $P\text{-value}$ estão acima de 0.01, o que indicaria uma distribuição equilibrada, 
            O resultado final do teste de frequência (monobit) para o Random mostra que ele obteve um $P\text{-value}$ de 0.275709 um valor mais proximo ao 0.0001.
            Por mais que seu valor seja inferior ao do SegureRandom, o Random ainda assim passou no teste de frequência (monobit) com uma distribuição de $P\text{-value}$ que é considerada boa.
            Um falso positivo, devido a maneira como o Random funciona, utilizando o LCG que é matematicamente desenhado para ter uma distribuição uniforme em complexidades simples.
        %PASSO 2 - Explicacao do resultado do (Teste de Frequência dentro de um bloco)
        \subsection{Teste de Frequência em Bloco}
            O teste de frequência em bloco é uma extensão do teste de frequência (monobit) que avalia a distribuição de bits '1' e '0' dentro de blocos específicos de uma sequência. Ele divide a sequência em blocos de comprimento M. E tenta determinar se o comprimento M é aproximadamente M/2.
            Caso o teste possua muitos zeros ou muitos uns, o teste passaria no teste de monibt, mas falharia no teste de frequência em bloco. Pois os blocos isolados estariam desequilibrados.
            
            \begin{figure}[htb]
                    \centering
                    \caption{Resultado Teste de Frequência em Bloco para SegureRandom}\label{fig:02_segure}
                    \includegraphics[width=0.8\textwidth]{NIST/02/02_segure.png}
                    \legend{Fonte: Elaborado pelo autor (2026)}
            \end{figure}

            Explicação dos valores:

            \begin{figure}[htb]
                    \centering
                    \caption{Resultado final do Teste de Frequência em Bloco para SegureRandom}\label{fig:02_final_segure}
                    \includegraphics[width=0.8\textwidth]{NIST/02/02_final_segure.png}
                    \legend{Fonte: Elaborado pelo autor (2026)}
            \end{figure}

            Explicação dos valores:

            \begin{figure}[htb]
                    \centering
                    \caption{Resultado Teste de Frequência em Bloco para Random}\label{fig:02_insegure}
                    \includegraphics[width=0.8\textwidth]{NIST/02/02_insegure.png}
                    \legend{Fonte: Elaborado pelo autor (2026)}
            \end{figure}

            Explicação dos valores:

            \begin{figure}[htb]
                    \centering
                    \caption{Resultado final do Teste de Frequência em Bloco para Random}\label{fig:02_final_insegure}
                    \includegraphics[width=0.8\textwidth]{NIST/02/02_final_insegure.png}
                    \legend{Fonte: Elaborado pelo autor (2026)}
            \end{figure}
            
        %PASSO 3 - Explicacao do resultado do (Teste de Sequencias)
        \subsection{Teste de Sequencias}
            Muitos textos explicativos

            \begin{figure}[htb]
                    \centering
                    \caption{Resultado Teste de Sequencia SegureRandom}\label{fig:03_segure}
                    \includegraphics[width=0.8\textwidth]{NIST/03/03_segure.png}
                    \legend{Fonte: Elaborado pelo autor (2026)}
            \end{figure}

            Explicação dos valores:

            \begin{figure}[htb]
                    \centering
                    \caption{Resultado Final do Teste de Sequencia SegureRandom}\label{fig:03_final_segure}
                    \includegraphics[width=0.8\textwidth]{NIST/03/03_final_segure.png}
                    \legend{Fonte: Elaborado pelo autor (2026)}
            \end{figure}

            Explicação dos valores:

            \begin{figure}[htb]
                    \centering
                    \caption{Resultado Teste de Sequencia Random}\label{fig:03_insegure}
                    \includegraphics[width=0.8\textwidth]{NIST/03/03_insegure.png}
                    \legend{Fonte: Elaborado pelo autor (2026)}
            \end{figure}

            Explicação dos valores:

            \begin{figure}[htb]
                    \centering
                    \caption{Resultado Final do Teste de Sequencia Random}\label{fig:03_final_insegure}
                    \includegraphics[width=0.8\textwidth]{NIST/03/03_final_insegure.png}
                    \legend{Fonte: Elaborado pelo autor (2026)}
            \end{figure}

            Explicação dos valores:

        %PASSO 4 - Explicacao do resultado do (Testes para Maior Sequencia de Uns em um bloco)

        \subsection{Testes para Maior Sequencia de Uns em um bloco}
            Muitos textos explicativos

            \begin{figure}[htb]
                    \centering
                    \caption{Resultado Testes para Maior Sequencia de Uns em um bloco SegureRandom}\label{fig:04_segure}
                    \includegraphics[width=0.8\textwidth]{NIST/04/04_segure.png}
                    \legend{Fonte: Elaborado pelo autor (2026)}
            \end{figure}

            Explicação dos valores:

            \begin{figure}[htb]
                    \centering
                    \caption{Resultado Final do Testes para Maior Sequencia de Uns em um bloco SegureRandom}\label{fig:04_final_segure}
                    \includegraphics[width=0.8\textwidth]{NIST/04/04_final_segure.png}
                    \legend{Fonte: Elaborado pelo autor (2026)}
            \end{figure}

            Explicação dos valores:

            \begin{figure}[htb]
                    \centering
                    \caption{Resultado Testes para Maior Sequencia de Uns em um bloco Random}\label{fig:04_insegure}
                    \includegraphics[width=0.8\textwidth]{NIST/04/04_insegure.png}
                    \legend{Fonte: Elaborado pelo autor (2026)}
            \end{figure}

            Explicação dos valores:

            \begin{figure}[htb]
                    \centering
                    \caption{Resultado Final do Testes para Maior Sequencia de Uns em um bloco Random}\label{fig:04_final_insegure}
                    \includegraphics[width=0.8\textwidth]{NIST/04/04_final_insegure.png}
                    \legend{Fonte: Elaborado pelo autor (2026)}
            \end{figure}

            Explicação dos valores:
        %PASSO 5 - Explicacao do resultado do (Testes de Posto de Matriz Binária)
        \subsection{Testes de Posto de Matriz Binária}
            Muitos textos explicativos

            \begin{figure}[htb]
                    \centering
                    \caption{Resultado Testes de Posto de Matriz Binária SegureRandom}\label{fig:05_segure}
                    \includegraphics[width=0.8\textwidth]{NIST/05/05_segure.png}
                    \legend{Fonte: Elaborado pelo autor (2026)}
            \end{figure}

            Explicação dos valores:

            \begin{figure}[htb]
                    \centering
                    \caption{Resultado Final do Testes de Posto de Matriz Binária SegureRandom}\label{fig:05_final_segure}
                    \includegraphics[width=0.8\textwidth]{NIST/05/05_final_segure.png}
                    \legend{Fonte: Elaborado pelo autor (2026)}
            \end{figure}

            Explicação dos valores:

            \begin{figure}[htb]
                    \centering
                    \caption{Resultado Testes de Posto de Matriz Binária Random}\label{fig:05_insegure}
                    \includegraphics[width=0.8\textwidth]{NIST/05/05_insegure.png}
                    \legend{Fonte: Elaborado pelo autor (2026)}
            \end{figure}

            Explicação dos valores:

            \begin{figure}[htb]
                    \centering
                    \caption{Resultado Final do Testes de Posto de Matriz Binária Random}\label{fig:05_final_insegure}
                    \includegraphics[width=0.8\textwidth]{NIST/05/05_final_insegure.png}
                    \legend{Fonte: Elaborado pelo autor (2026)}
            \end{figure}

            Explicação dos valores:
        %PASSO 6 - Explicacao do resultado do (Teste da Transformada Discreta de Fourier (Espectral))
        \subsection{Teste da Transformada Discreta de Fourier (Espectral)}
            Muitos textos explicativos

            \begin{figure}[htb]
                    \centering
                    \caption{Resultado Teste da Transformada Discreta de Fourier (Espectral) SegureRandom}\label{fig:06_segure}
                    \includegraphics[width=0.8\textwidth]{NIST/06/06_segure.png}
                    \legend{Fonte: Elaborado pelo autor (2026)}
            \end{figure}

            Explicação dos valores:

            \begin{figure}[htb]
                    \centering
                    \caption{Resultado Final do Teste da Transformada Discreta de Fourier (Espectral) SegureRandom}\label{fig:06_final_segure}
                    \includegraphics[width=0.8\textwidth]{NIST/06/06_final_segure.png}
                    \legend{Fonte: Elaborado pelo autor (2026)}
            \end{figure}

            Explicação dos valores:

            \begin{figure}[htb]
                    \centering
                    \caption{Resultado Teste da Transformada Discreta de Fourier (Espectral) Random}\label{fig:06_insegure}
                    \includegraphics[width=0.8\textwidth]{NIST/06/06_insegure.png}
                    \legend{Fonte: Elaborado pelo autor (2026)}
            \end{figure}

            Explicação dos valores:

            \begin{figure}[htb]
                    \centering
                    \caption{Resultado Final do Teste da Transformada Discreta de Fourier (Espectral) Random}\label{fig:06_final_insegure}
                    \includegraphics[width=0.8\textwidth]{NIST/06/06_final_insegure.png}
                    \legend{Fonte: Elaborado pelo autor (2026)}
            \end{figure}

            Explicação dos valores:
        %PASSO 7 - Explicacao do resultado do (Teste de Correspondência de Modelos não Sobrepostos)
        \subsection{Teste de Correspondência de Modelos não Sobrepostos}
            Muitos textos explicativos

            \begin{figure}[htb]
                    \centering
                    \caption{Resultado Final Teste de Correspondência de Modelos não Sobrepostos SegureRandom}\label{fig:07_final_segure_pt1}
                    \includegraphics[width=0.8\textwidth]{NIST/07/07_final_segure_pt1.png}
                    \legend{Fonte: Elaborado pelo autor (2026)}
            \end{figure}

            Explicação dos valores:

            \begin{figure}[htb]
                    \centering
                    \caption{Resultado Final do Teste de Correspondência de Modelos não Sobrepostos SegureRandom}\label{fig:07_final_segure_pt2}
                    \includegraphics[width=0.8\textwidth]{NIST/07/07_final_segure_pt2.png}
                    \legend{Fonte: Elaborado pelo autor (2026)}
            \end{figure}

            Explicação dos valores:

            \begin{figure}[htb]
                    \centering
                    \caption{Resultado Final Teste de Correspondência de Modelos não Sobrepostos Random}\label{fig:07_final_insegure_pt1}
                    \includegraphics[width=0.8\textwidth]{NIST/07/07_final_insegure_pt1.png}
                    \legend{Fonte: Elaborado pelo autor (2026)}
            \end{figure}

            Explicação dos valores:

            \begin{figure}[htb]
                    \centering
                    \caption{Resultado Final do Teste de Correspondência de Modelos não Sobrepostos Random}\label{fig:07_final_insegure_pt2}
                    \includegraphics[width=0.8\textwidth]{NIST/07/07_final_insegure_pt2.png}
                    \legend{Fonte: Elaborado pelo autor (2026)}
            \end{figure}

            Explicação dos valores:
        %PASSO 8 - Explicacao do resultado do (Teste de Correspondência de Modelos Sobrepostos)
        \subsection{Teste de Correspondência de Modelos Sobrepostos}
            Muitos textos explicativos

            \begin{figure}[htb]
                    \centering
                    \caption{Resultado Teste de Correspondência de Modelos Sobrepostos SegureRandom}\label{fig:08_segure}
                    \includegraphics[width=0.8\textwidth]{NIST/08/08_segure.png}
                    \legend{Fonte: Elaborado pelo autor (2026)}
            \end{figure}

            Explicação dos valores:

            \begin{figure}[htb]
                    \centering
                    \caption{Resultado Final do Teste de Correspondência de Modelos Sobrepostos SegureRandom}\label{fig:08_final_segure}
                    \includegraphics[width=0.8\textwidth]{NIST/08/08_final_segure.png}
                    \legend{Fonte: Elaborado pelo autor (2026)}
            \end{figure}

            Explicação dos valores:

            \begin{figure}[htb]
                    \centering
                    \caption{Resultado Teste de Correspondência de Modelos Sobrepostos Random}\label{fig:08_insegure}
                    \includegraphics[width=0.8\textwidth]{NIST/08/08_insegure.png}
                    \legend{Fonte: Elaborado pelo autor (2026)}
            \end{figure}

            Explicação dos valores:

            \begin{figure}[htb]
                    \centering
                    \caption{Resultado Final do Teste de Correspondência de Modelos Sobrepostos Random}\label{fig:08_final_insegure}
                    \includegraphics[width=0.8\textwidth]{NIST/08/08_final_insegure.png}
                    \legend{Fonte: Elaborado pelo autor (2026)}
            \end{figure}

            Explicação dos valores:
        %PASSO 9 - Explicacao do resultado do (Teste "Estatístico Universal" de Maurer)
        \subsection{Teste `Estatístico Universal'' de Maurer}
            Muitos textos explicativos

            \begin{figure}[htb]
                    \centering
                    \caption{Resultado Teste `Estatístico Universal'' de Maurer SegureRandom}\label{fig:09_segure}
                    \includegraphics[width=0.8\textwidth]{NIST/09/09_segure.png}
                    \legend{Fonte: Elaborado pelo autor (2026)}
            \end{figure}

            Explicação dos valores:

            \begin{figure}[htb]
                    \centering
                    \caption{Resultado Final do Teste `Estatístico Universal'' de Maurer SegureRandom}\label{fig:09_final_segure}
                    \includegraphics[width=0.8\textwidth]{NIST/09/09_final_segure.png}
                    \legend{Fonte: Elaborado pelo autor (2026)}
            \end{figure}

            Explicação dos valores:

            \begin{figure}[htb]
                    \centering
                    \caption{Resultado Teste `Estatístico Universal'' de Maurer Random}\label{fig:09_insegure}
                    \includegraphics[width=0.8\textwidth]{NIST/09/09_insegure.png}
                    \legend{Fonte: Elaborado pelo autor (2026)}
            \end{figure}

            Explicação dos valores:

            \begin{figure}[htb]
                    \centering
                    \caption{Resultado Final do Teste `Estatístico Universal'' de Maurer Random}\label{fig:09_final_insegure}
                    \includegraphics[width=0.8\textwidth]{NIST/09/09_final_insegure.png}

                    \legend{Fonte: Elaborado pelo autor (2026)}
            \end{figure}

            Explicação dos valores:

            Explicação dos valores:
        %PASSO 10 - Explicacao do resultado do (Teste de Complexidade Linear)
        \subsection{Teste de Complexidade Linear}
            Muitos textos explicativos

            \begin{figure}[htb]
                    \centering
                    \caption{Resultado Teste de Complexidade Linear SegureRandom}\label{fig:10_segure}
                    \includegraphics[width=0.8\textwidth]{NIST/10/10_segure.png}
                    \legend{Fonte: Elaborado pelo autor (2026)}
            \end{figure}

            Explicação dos valores:

            \begin{figure}[htb]
                    \centering
                    \caption{Resultado Final do Teste de Complexidade Linear SegureRandom}\label{fig:10_final_segure}
                    \includegraphics[width=0.8\textwidth]{NIST/10/10_final_segure.png}
                    \legend{Fonte: Elaborado pelo autor (2026)}
            \end{figure}

            Explicação dos valores:

            \begin{figure}[htb]
                    \centering
                    \caption{Resultado Teste de Complexidade Linear Random}\label{fig:10_insegure}
                    \includegraphics[width=0.8\textwidth]{NIST/10/10_insegure.png}
                    \legend{Fonte: Elaborado pelo autor (2026)}
            \end{figure}

            Explicação dos valores:

            \begin{figure}[htb]
                    \centering
                    \caption{Resultado Final do Teste de Complexidade Linear Random}\label{fig:10_final_insegure}
                    \includegraphics[width=0.8\textwidth]{NIST/10/10_final_insegure.png}
                    \legend{Fonte: Elaborado pelo autor (2026)}
            \end{figure}

            Explicação dos valores:
        %PASSO 11 - Explicacao do resultado do (Teste Serial)
        \subsection{Teste Serial}
            Muitos textos explicativos

            \begin{figure}[htb]
                    \centering
                    \caption{Resultado Teste Serial SegureRandom}\label{fig:11_segure}
                    \includegraphics[width=0.8\textwidth]{NIST/11/11_segure.png}
                    \legend{Fonte: Elaborado pelo autor (2026)}
            \end{figure}

            Explicação dos valores:

            \begin{figure}[htb]
                    \centering
                    \caption{Resultado Final do Teste Serial SegureRandom}\label{fig:11_final_segure}
                    \includegraphics[width=0.8\textwidth]{NIST/11/11_final_segure.png}
                    \legend{Fonte: Elaborado pelo autor (2026)}
            \end{figure}

            Explicação dos valores:

            \begin{figure}[htb]
                    \centering
                    \caption{Resultado Teste Serial Random}\label{fig:11_insegure}
                    \includegraphics[width=0.8\textwidth]{NIST/11/11_insegure.png}
                    \legend{Fonte: Elaborado pelo autor (2026)}
            \end{figure}

            Explicação dos valores:

            \begin{figure}[htb]
                    \centering
                    \caption{Resultado Final do Teste Serial Random}\label{fig:11_final_insegure}
                    \includegraphics[width=0.8\textwidth]{NIST/11/11_final_insegure.png}
                    \legend{Fonte: Elaborado pelo autor (2026)}
            \end{figure}

            Explicação dos valores:
        %PASSO 12 - Explicacao do resultado do (Teste de Entropia Aproximada)
        \subsection{Teste de Entropia Aproximada}
            Muitos textos explicativos

            \begin{figure}[htb]
                    \centering
                    \caption{Resultado Teste de Entropia Aproximada SegureRandom}\label{fig:12_segure}
                    \includegraphics[width=0.8\textwidth]{NIST/12/12_segure.png}
                    \legend{Fonte: Elaborado pelo autor (2026)}
            \end{figure}

            Explicação dos valores:

            \begin{figure}[htb]
                    \centering
                    \caption{Resultado Final do Teste de Entropia Aproximada SegureRandom}\label{fig:12_final_segure}
                    \includegraphics[width=0.8\textwidth]{NIST/12/12_final_segure.png}
                    \legend{Fonte: Elaborado pelo autor (2026)}
            \end{figure}

            Explicação dos valores:

            \begin{figure}[htb]
                    \centering
                    \caption{Resultado Teste de Entropia Aproximada Random}\label{fig:12_insegure}
                    \includegraphics[width=0.8\textwidth]{NIST/12/12_insegure.png}
                    \legend{Fonte: Elaborado pelo autor (2026)}
            \end{figure}

            Explicação dos valores:

            \begin{figure}[htb]
                    \centering
                    \caption{Resultado Final do Teste de Entropia Aproximada Random}\label{fig:12_final_insegure}
                    \includegraphics[width=0.8\textwidth]{NIST/12/12_final_insegure.png}
                    \legend{Fonte: Elaborado pelo autor (2026)}
            \end{figure}

            Explicação dos valores:
        %PASSO 13 - Explicacao do resultado do (Teste de Somas Cumulativas (Cusums))
        \subsection{Teste de Somas Cumulativas (Cusums)}
            Muitos textos explicativos

            \begin{figure}[htb]
                    \centering
                    \caption{Resultado Teste de Somas Cumulativas (Cusums) SegureRandom}\label{fig:13_segure}
                    \includegraphics[width=0.8\textwidth]{NIST/13/13_segure.png}
                    \legend{Fonte: Elaborado pelo autor (2026)}
            \end{figure}

            Explicação dos valores:

            \begin{figure}[htb]
                    \centering
                    \caption{Resultado Final do Teste de Somas Cumulativas (Cusums) SegureRandom}\label{fig:13_final_segure}
                    \includegraphics[width=0.8\textwidth]{NIST/13/13_final_segure.png}
                    \legend{Fonte: Elaborado pelo autor (2026)}
            \end{figure}

            Explicação dos valores:

            \begin{figure}[htb]
                    \centering
                    \caption{Resultado Teste de Somas Cumulativas (Cusums) Random}\label{fig:13_insegure}
                    \includegraphics[width=0.8\textwidth]{NIST/13/13_insegure.png}
                    \legend{Fonte: Elaborado pelo autor (2026)}
            \end{figure}

            Explicação dos valores:

            \begin{figure}[htb]
                    \centering
                    \caption{Resultado Final do Teste de Somas Cumulativas (Cusums) Random}\label{fig:13_final_insegure}
                    \includegraphics[width=0.8\textwidth]{NIST/13/13_final_insegure.png}
                    \legend{Fonte: Elaborado pelo autor (2026)}
            \end{figure}

            Explicação dos valores:
        %PASSO 14 - Explicacao do resultado do (Teste de Excursões Aleatórias)
        \subsection{Teste de Excursões Aleatórias}
            Muitos textos explicativos

            \begin{figure}[htb]
                    \centering
                    \caption{Resultado Final do Teste de Excursões Aleatórias SegureRandom}\label{fig:14_final_segure}
                    \includegraphics[width=0.8\textwidth]{NIST/14/14_final_segure.png}
                    \legend{Fonte: Elaborado pelo autor (2026)}
            \end{figure}

            Explicação dos valores:

            \begin{figure}[htb]
                    \centering
                    \caption{Resultado Final do Teste de Excursões Aleatórias Random}\label{fig:14_final_insegure}
                    \includegraphics[width=0.8\textwidth]{NIST/14/14_final_insegure.png}
                    \legend{Fonte: Elaborado pelo autor (2026)}
            \end{figure}

            Explicação dos valores:
        %PASSO 15 - Explicacao do resultado do (Teste de Variante de Excursões Aleatórias)
        \subsection{Teste de Variante de Excursões Aleatórias}
            Muitos textos explicativos

            \begin{figure}[htb]
                    \centering
                    \caption{Resultado Final do Teste de Variante de Excursões Aleatórias SegureRandom}\label{fig:15_final_segure}
                    \includegraphics[width=0.8\textwidth]{NIST/15/15_final_segure.png}
                    \legend{Fonte: Elaborado pelo autor (2026)}
            \end{figure}

            Explicação dos valores:


            \begin{figure}[htb]
                    \centering
                    \caption{Resultado Final do Teste de Variante de Excursões Aleatórias Random}\label{fig:15_final_insegure}
                    \includegraphics[width=0.8\textwidth]{NIST/15/15_final_insegure.png}
                    \legend{Fonte: Elaborado pelo autor (2026)}
            \end{figure}

            Explicação dos valores:
        %Passo 16 Explicacao do porque esse teste é somente uma das etapas pra considerar um gerador seguro e não a fonte da verdade

\chapter{Coisas possivelmente Uteis em algum momento}
    PRNG é uma sigla para “Pseudo Random Number Generator” ou “Gerador de Números Pseudo Aleatórios”. 
    Seu significado é que a sua geração de números não segue uma estrutura aleatória. 
    Ele utiliza um ponto inicial denominado “seed” para dar início a sua geração de números. 
    A seed é um valor base para ser utilizado no PRNG. 
    Utilizando a mesma seed, o resultado final do PRNG será o mesmo.

    O PRNG funciona a partir de uma função que recebe a seed por parâmetro. 
    E através de um cálculo, ele retorna o resultado dessa fórmula. 
    Uma das fórmulas mais comuns para esse tipo de gerador de pseudo números aleatórios é no modelo LCG “Gerador Congruencial Linear” 

    \begin{equation}
        x_{i+1} = (A \cdot x_i + C) \pmod{M}
        \label{eq:lca-appendix}
    \end{equation}

``Onde o valor inicial "x0" é a seed, A e C são parâmetros inteiros que podem ter um valor já pré-definido ou passados diretamente, M é o espaço dos valores aleatórios possíveis e "mod" é o restante da divisão inteira.''


    \begin{figure}[htb]
        \centering
        \caption{Modelo simples de um PRNG com LCG}\label{fig:model_prng_lcg_appendix}
        \includegraphics[width=0.8\textwidth]{Modelos/model_prng_lcg.png}
        \legend{Fonte: Elaborado pelo autor (2026)}
    \end{figure}

    Mesmo após inúmeras tentativas. Devido a seed permanecer inalterada. 
    O resultado permanece o mesmo em todos os resultados obtidos.

    Por mais que tenha se passado um valor fixo na seed no exemplo prático, a seed é definida pelo sistema operacional, esse número é denominado entropia. 
    Que se tem por referência a quantidade de imprecisão de um sistema. 
    Quanto maior o número da entropia, maior é a incerteza. 
    Esse valor de entropia é gerado a partir de certos eventos que ocorrem no momento de sua captura, como por exemplo o movimento do mouse, os inputs no teclado etc. 
    Essas capturas são convertidas em um valor numérico para ser usado na seed.

    Esses inputs que o sistema captura podem ser previsíveis, tornando a entropia causada menor. 
    Sendo possível pré-determinar o valor gerado ao final para seed. 
    Apesar das aplicações para o dia a dia ou em pequenos projetos serem o suficiente. 
    Esse método  não é seguro para aplicações de alto risco.


    \begin{figure}[htb]
        \centering
        \caption{Resultado do script LCG}\label{fig:result_prng_lcg_appendix}
        \includegraphics[width=0.8\textwidth]{Resultados/result_prng_lcg.png}
        \legend{Fonte: Elaborado pelo autor (2026)}
    \end{figure}

